\chapter{PHƯƠNG PHÁP NGHIÊN CỨU VÀ THIẾT KẾ MÔ HÌNH NEUMF}
\label{chap:thietke_trienkhai}

\section{Kiến trúc tổng thể của hệ thống nghiên cứu}
\label{sec:system_architecture}

Trước khi trình bày chi tiết từng thành phần, mục này mô tả kiến trúc tổng thể của toàn bộ hệ thống thực nghiệm -- từ dữ liệu thô đến kết quả đánh giá cuối cùng -- nhằm cung cấp một bức tranh toàn cảnh về luồng xử lý và mối liên hệ giữa các thành phần. Hệ thống được tổ chức thành sáu tầng xử lý tuần tự (Hình~\ref{fig:system_architecture}): Dữ liệu nguồn, Tiền xử lý, Phân chia dữ liệu \& Lấy mẫu, Mô hình hóa, Huấn luyện \& Tinh chỉnh và Đánh giá, cùng một nhánh Đối chứng (Baseline) vận hành song song.

\begin{figure}[htbp]
    \centering
    \resizebox{\textwidth}{!}{%
    \begin{tikzpicture}[
        node distance=0.9cm,
        stage/.style={rectangle, draw=black!60, very thick, rounded corners=2pt, text width=13.5cm, align=left, font=\small, inner sep=8pt},
        side/.style={rectangle, draw=black!60, very thick, rounded corners=2pt, text width=4.6cm, align=left, font=\small, inner sep=6pt, fill=gray!8},
        arrow/.style={->, >=stealth, thick, color=black!70}
    ]
        \node[stage, fill=gray!12] (data) at (0,0)
            {\textbf{TẦNG 1 -- DỮ LIỆU NGUỒN}\\[2pt]
             H\&M Personalized Fashion Recommendations \cite{hm2022kaggle}: 31,8 triệu giao dịch (2018-09-20 $\rightarrow$ 2020-09-22) $\rightarrow$
             lấy mẫu \textbf{hm500k} theo khách hàng (500.269 giao dịch, 21.599 khách, giữ trọn lịch sử 2 năm). Chỉ dùng customer\_id, article\_id, t\_dat.};

        \node[stage, fill=blue!8, below=of data] (prep)
            {\textbf{TẦNG 2 -- TIỀN XỬ LÝ DỮ LIỆU}\\[2pt]
             (1) Gộp giao dịch lặp thành cặp (User, Item) duy nhất $\rightarrow$
             (2) Lọc $k$-core ($k=10$) lặp đến hội tụ $\rightarrow$
             (3) Re-indexing $u \in \{0,\dots,M-1\}$, $i \in \{0,\dots,N-1\}$ $\rightarrow$
             (4) Ma trận phản hồi ẩn nhị phân $Y \in \{0,1\}^{M\times N}$ (7.519 $\times$ 10.345).};

        \node[stage, fill=blue!8, below=of prep] (split)
            {\textbf{TẦNG 3 -- PHÂN CHIA DỮ LIỆU \& LẤY MẪU}\\[2pt]
             Chia theo \textbf{một mốc thời gian chung}: Train $<$ 01/07/2020 $\le$ Validation $<$ 29/07/2020 $\le$ Test (giao thức chính);
             Leave-One-Out theo thời gian (giao thức phụ) $\rightarrow$ Kiểm tra Disjoint bắt buộc (Công thức~\eqref{eq:disjoint_assert}) $\rightarrow$
             Negative Sampling động chỉ từ tập huấn luyện.};

        \node[stage, fill=green!8, below=of split] (model)
            {\textbf{TẦNG 4 -- MÔ HÌNH HÓA}\\[2pt]
             \textbf{GMF} (tuyến tính, Công thức~\eqref{eq:gmf_predict}) và \textbf{MLP} (phi tuyến, tháp $[2d{\to}d{\to}d/2{\to}d/4]$) $\rightarrow$
             hợp nhất thành \textbf{NeuMF} ở lớp gần cuối (Early Fusion); hai biến thể NeuMF-Scratch (khởi tạo ngẫu nhiên) và NeuMF-Pretrained (khởi tạo từ GMF, MLP đã huấn luyện).
             Mở rộng: \textbf{Late Fusion} -- trộn điểm GMF + MLP (và BPR-MF + MLP) huấn luyện riêng.};

        \node[stage, fill=orange!10, below=of model] (train)
            {\textbf{TẦNG 5 -- HUẤN LUYỆN \& TINH CHỈNH (chỉ dùng Validation)}\\[2pt]
             Tối thiểu hóa $\mathcal{L}_{BCE}$ bằng Adam, Early Stopping theo NDCG@10 trên Validation $\rightarrow$
             Tuning cùng ngân sách cho mọi mô hình (6 cấu hình/mô hình), ghi toàn bộ vào nhật ký tuning $\rightarrow$ chọn cấu hình tốt nhất.};

        \node[stage, fill=red!8, below=of train] (eval)
            {\textbf{TẦNG 6 -- ĐÁNH GIÁ CUỐI TRÊN TEST (khoá tập test)}\\[2pt]
             Huấn luyện lại trên Train $\cup$ Validation đúng số epoch tốt nhất $\rightarrow$
             Full Ranking, 3 seed $\rightarrow$ NDCG@10 (chính), Recall@10, HR@10, Precision@10, NDCG@5 $\rightarrow$
             kiểm định cặp theo user (Wilcoxon, bootstrap, Holm) $\rightarrow$ phân tích phụ: Long-tail, Cold/Warm, Beyond-accuracy, Sampled-99, độ trễ.};

        \node[side, text width=5.0cm, right=1.1cm of train] (baselinebox)
            {\textbf{ĐỐI CHỨNG (BASELINES)}\\[2pt]
             Random, Most Popular\\
             BPR-MF (mục~\ref{subsec:bpr}), tune cùng ngân sách\\[2pt]
             \textit{Mở rộng:} ItemKNN, UserKNN (mục~\ref{subsec:memory_model_based})};

        \draw[arrow] (data) -- (prep);
        \draw[arrow] (prep) -- (split);
        \draw[arrow] (split) -- (model);
        \draw[arrow] (model) -- (train);
        \draw[arrow] (train) -- (eval);
        \draw[arrow] (baselinebox.south) |- (eval.east);
    \end{tikzpicture}}
    \caption{Kiến trúc tổng thể của hệ thống nghiên cứu, từ dữ liệu thô H\&M đến đánh giá cuối trên tập test}
    \label{fig:system_architecture}
\end{figure}

Ba đặc điểm thiết kế xuyên suốt cần nhấn mạnh: \textit{(i)} tầng tiền xử lý và tầng phân chia dữ liệu tách biệt khỏi tầng mô hình hóa, nên mọi mô hình -- NeuMF lẫn baseline -- dùng đúng một bộ dữ liệu, một tập ứng viên và một quy tắc loại sản phẩm đã mua, đảm bảo so sánh công bằng; \textit{(ii)} mọi quyết định chọn mô hình và siêu tham số chỉ dựa trên tập xác thực, tập kiểm thử chỉ được dùng ở tầng 6, qua kịch bản có khoá và ghi nhật ký mọi lần truy cập (mục~\ref{sec:protocol}); \textit{(iii)} baseline BPR-MF được tinh chỉnh với cùng ngân sách như NeuMF, vì các nghiên cứu đối chứng gần đây cho thấy lợi thế của mô hình học sâu thường biến mất khi baseline được tune đúng cách \parencite{rendle2020ncfmf, dacrema2019progress}.

\section{Phân tích và tiền xử lý dữ liệu H\&M Personalized Fashion Recommendations}
\label{sec:hm_pipeline}

Bộ dữ liệu thực nghiệm của dự án là H\&M Personalized Fashion Recommendations \parencite{hm2022kaggle}, phản ánh hoạt động mua sắm thời trang thực tế của một chuỗi bán lẻ lớn. Quy mô của bộ dữ liệu đặt ra thách thức về chi phí tính toán mà dự án phải giải quyết trước khi áp dụng quy trình xử lý chuẩn.

\subsection{Đặc tả bộ dữ liệu và thách thức về quy mô}
\label{subsec:hm_description}

Tệp giao dịch gốc (\texttt{transactions\_train.csv}, 3,49 GB) chứa 31.788.324 dòng giao dịch mua hàng thời trang từ 20/09/2018 đến 22/09/2020, ghi nhận 1.362.281 khách hàng và 104.547 sản phẩm riêng biệt. Ba trường được dùng để xây dựng sự kiện tương tác:
\begin{itemize}
    \item \textbf{customer\_id:} Mã định danh khách hàng (User ID), dạng chuỗi mã băm 64 ký tự.
    \item \textbf{article\_id:} Mã định danh sản phẩm (Item ID), mã số 10 chữ số -- phải đọc dưới dạng chuỗi và giữ số 0 ở đầu.
    \item \textbf{t\_dat:} Ngày mua, dùng để sắp xếp tương tác theo thời gian (chỉ ở mức ngày, không có giờ).
\end{itemize}
Trường giá (\texttt{price}) và hai tệp thuộc tính \texttt{articles.csv}, \texttt{customers.csv} \emph{không} được dùng làm đầu vào của bất kỳ mô hình nào: đề tài giới hạn ở lọc cộng tác thuần dựa trên ID (mục~\ref{subsec:phamvi}); hai tệp thuộc tính chỉ dùng để hiển thị tên sản phẩm trong hệ thống demo (mục~\ref{sec:demo_design}). Ràng buộc này được kiểm chứng tự động: bộ kiểm thử của dự án khẳng định bộ đọc dữ liệu chỉ mở tệp giao dịch.

Đánh giá Full Ranking (mục~\ref{sec:metrics}) phải chấm điểm \emph{mọi} sản phẩm cho mỗi người dùng kiểm thử, nên chi phí tỉ lệ với tích số người dùng và số sản phẩm. Bảng~\ref{tab:kcore_hm_full} cho thấy ở quy mô đầy đủ, ngay cả với $k=10$ vẫn còn 633.130 người dùng và 80.265 sản phẩm -- tức hàng chục tỷ cặp cần chấm cho mỗi mô hình, vượt xa khả năng của phần cứng sẵn có (GPU 4 GB, RAM 16 GB).

\begin{table}[htbp]
    \centering
    \caption{Độ nhạy $k$-core trên toàn bộ dữ liệu H\&M (31.788.324 dòng, quy mô đầy đủ)}
    \label{tab:kcore_hm_full}
    \begin{tabular}{|c|c|c|c|c|}
        \hline
        \textbf{$k$} & \textbf{Số người dùng} & \textbf{Số sản phẩm} & \textbf{Số tương tác} & \textbf{Mật độ (\%)} \\
        \hline
        3 & 1.074.388 & 96.264 & 26.873.935 & 0,0260 \\
        \hline
        5 & 889.062 & 90.690 & 26.215.294 & 0,0325 \\
        \hline
        10 & 633.130 & 80.265 & 24.426.258 & 0,0481 \\
        \hline
    \end{tabular}
\end{table}

\subsection{Lấy mẫu hm500k theo khách hàng}
\label{subsec:hm_sampling}

Để giữ Full Ranking khả thi mà vẫn bảo toàn tính thời gian của dữ liệu, dự án lấy một mẫu con \emph{theo khách hàng} thay vì theo dòng. Lý do: lấy mẫu theo dòng (hoặc cắt $n$ dòng đầu) sẽ cắt vụn lịch sử của từng khách và -- vì tệp gốc sắp theo ngày -- chỉ phủ vài ngày giao dịch, làm hỏng cả bước lọc $k$-core lẫn việc chia dữ liệu theo thời gian. Quy trình lấy mẫu gồm hai lượt đọc tệp theo lô (giữ bộ nhớ đỉnh khoảng 0,5 GB):
\begin{enumerate}
    \item \textbf{Lượt 1 -- chọn khách:} Băm \texttt{customer\_id} bằng một khoá cố định vào 100.000 nhóm (hàm băm phân khách vào nhóm một cách giả ngẫu nhiên, độc lập với hành vi mua); đếm số giao dịch của từng nhóm; lấy lần lượt các nhóm theo chỉ số cho tới khi tổng số giao dịch vừa vượt 500.000. Kết quả: 1.597 nhóm (1,60\% khách hàng).
    \item \textbf{Lượt 2 -- ghi mẫu:} Giữ \emph{toàn bộ} giao dịch của các khách được chọn, trải đủ 20/09/2018 -- 22/09/2020, giữ nguyên các cột như tệp gốc.
\end{enumerate}
Mẫu thu được -- gọi là \textbf{hm500k} -- gồm 500.269 giao dịch của 21.599 khách hàng trên 59.483 sản phẩm. Vì dựa trên hàm băm với khoá cố định, chạy lại luôn cho đúng cùng một tệp (tái lập được), và toàn bộ quá trình mất khoảng 70 giây.

\subsection{Xây dựng ma trận tương tác phản hồi ẩn}
\label{subsec:interaction_matrix}

Một khách có thể mua cùng một sản phẩm nhiều lần; do đó bước đầu tiên là gộp các giao dịch lặp thành một cặp User--Item duy nhất (giữ lại ngày mua đầu tiên, ngày mua cuối cùng và số lần mua) \emph{trước} khi lọc $k$-core. Thứ tự này quan trọng: nếu lọc $k$-core trực tiếp trên dòng giao dịch thô, bậc của một User/Item bị thổi phồng bởi các lượt mua lặp và tập sau lọc không phản ánh đúng số đối tác riêng biệt. Trên hm500k, 500.269 giao dịch gộp lại còn 429.964 cặp duy nhất.

Sau khi gộp, lọc $k$-core được áp dụng lặp đến khi hội tụ (loại một sản phẩm có thể khiến một người dùng rơi xuống dưới ngưỡng và ngược lại), để mỗi thực thể còn lại có ít nhất $k$ đối tác riêng biệt. Toàn bộ ID được tái lập chỉ số về số nguyên liên tục:
\begin{equation}
    u \in \{0, 1, \dots, M-1\}, \quad i \in \{0, 1, \dots, N-1\}
\end{equation}
và ma trận tương tác nhị phân $Y \in \{0, 1\}^{M \times N}$ được tạo với $y_{ui} = 1$ nếu User $u$ đã mua Item $i$. Phản hồi nhị phân là thiết lập duy nhất dùng cho kết quả; mọi thống kê phụ thuộc dữ liệu (độ phổ biến, nhóm head/tail, nhóm người dùng ít tương tác, tập mẫu âm) đều chỉ tính trên tập huấn luyện.

\subsection{Lựa chọn ngưỡng lọc \texorpdfstring{$k$}{k}-core}
\label{subsec:hm_apply}

Bảng~\ref{tab:kcore_hm} khảo sát độ nhạy của hm500k theo ngưỡng $k$. Ngưỡng chuẩn trong tài liệu là $k=5$, nhưng khi đó catalog còn 23.067 sản phẩm và đánh giá Full Ranking ước tính cần khoảng 10,7 GB RAM -- vượt giới hạn của máy thực nghiệm (16 GB, đã tính hệ điều hành và dữ liệu huấn luyện). Dự án vì vậy chọn $\mathbf{k=10}$: 7.519 người dùng, 10.345 sản phẩm, 220.292 cặp tương tác, mật độ 0,283\%; RAM đỉnh đo thực tế khi huấn luyện và đánh giá một seed là 4,1 GB. Cái giá của lựa chọn này là dữ liệu thiên về khách mua nhiều -- được ghi nhận là một hạn chế (mục~\ref{sec:hanche}).

\begin{table}[htbp]
    \centering
    \caption{Độ nhạy của mẫu hm500k theo ngưỡng lọc $k$-core (429.964 cặp trước khi lọc)}
    \label{tab:kcore_hm}
    \begin{tabular}{|c|c|c|c|c|}
        \hline
        \textbf{$k$} & \textbf{Số người dùng} & \textbf{Số sản phẩm} & \textbf{Số tương tác} & \textbf{Mật độ (\%)} \\
        \hline
        3 & 16.377 & 33.620 & 387.733 & 0,070 \\
        \hline
        5 & 12.900 & 23.067 & 340.267 & 0,114 \\
        \hline
        \textbf{10} & \textbf{7.519} & \textbf{10.345} & \textbf{220.292} & \textbf{0,283} \\
        \hline
        15 & 3.017 & 3.186 & 83.088 & 0,864 \\
        \hline
    \end{tabular}
\end{table}

\subsection{Phân chia dữ liệu theo thời gian}
\label{subsec:split_sampling}

\paragraph{Giao thức chính -- một mốc thời gian chung.} Mỗi cặp (User, Item) được xếp theo \emph{ngày mua đầu tiên} và chia bằng hai mốc cố định cho mọi người dùng:
\begin{equation}
    \text{Train} < \text{01/07/2020} \le \text{Validation} < \text{29/07/2020} \le \text{Test} \;\text{(đến 22/09/2020)}
    \label{eq:global_split}
\end{equation}
Tập xác thực phủ 4 tuần, tập kiểm thử phủ 8 tuần. Tập xác thực chỉ giữ người dùng và sản phẩm đã xuất hiện trong tập huấn luyện; tập kiểm thử chỉ giữ người dùng và sản phẩm đã xuất hiện trong tập huấn luyện $\cup$ xác thực (mô hình lọc cộng tác thuần ID không thể chấm điểm cho thực thể chưa từng thấy). Mỗi người dùng có thể có \emph{nhiều} sản phẩm đúng, và sản phẩm người dùng đã mua trước đó không được tính là đích -- mô hình dự đoán những món \emph{mới} đối với người dùng. Cách chia này bảo đảm mô hình không bao giờ thấy tương tác xảy ra sau thời điểm cần dự đoán: khi tinh chỉnh, mô hình chỉ học từ dữ liệu trước 01/07/2020; khi chấm kiểm thử, chỉ học từ dữ liệu trước 29/07/2020 -- đúng với tình huống triển khai thật \parencite{meng2020splitting}.

\paragraph{Huấn luyện lại trước khi chấm kiểm thử.} Tập xác thực chỉ dùng để chọn cấu hình và số epoch. Trước khi chấm tập kiểm thử, mỗi mô hình được huấn luyện lại từ đầu trên tập huấn luyện $\cup$ xác thực (mọi cặp trước 29/07/2020) với cấu hình tốt nhất, chạy đúng số epoch tốt nhất đã chọn trên tập xác thực; lúc này không còn tập xác thực để dừng sớm, và tập kiểm thử không được dùng để dừng. Most Popular và BPR-MF cũng được khớp trên tập huấn luyện $\cup$ xác thực. Nhờ vậy, khi chấm kiểm thử mô hình đã thấy toàn bộ quá khứ như lúc triển khai thật, thay vì bỏ trống 4 tuần sát kỳ dự đoán -- một kiểm tra chỉ trên tập xác thực cho thấy bỏ trống 4 tuần như vậy làm NDCG@10 của Most Popular và BPR-MF giảm khoảng 17--19\%, lớn hơn chênh lệch giữa các mô hình. Bước này được bổ sung \emph{sau} lần chấm kiểm thử đầu tiên và được ghi nhận là lệch kế hoạch (mục~\ref{sec:protocol}).

\paragraph{Giao thức phụ -- Leave-One-Out theo thời gian.} Theo đúng NCF gốc \parencite{he2017ncf}: với mỗi người dùng, sản phẩm mua cuối cùng làm tập kiểm thử, sản phẩm áp chót làm tập xác thực, phần còn lại làm tập huấn luyện. Giao thức này được giữ để đối chiếu với tài liệu, \emph{không} dùng để chọn mô hình, vì nó rò rỉ tương lai: đo trên hm500k, tính trung bình theo người dùng, 11,2\% tương tác trong tập huấn luyện (của toàn bộ người dùng) xảy ra \emph{sau} ngày mua sản phẩm kiểm thử của người dùng đó -- tức mô hình được học từ những gì xảy ra sau thời điểm cần dự đoán; ngoài ra 60,9\% người dùng có sản phẩm xác thực và kiểm thử được mua cùng một ngày.

\paragraph{Kiểm tra không giao nhau.} Sau khi chia, hệ thống bắt buộc kiểm tra phần giao giữa từng cặp tập ở cấp độ cặp (User, Item):
\begin{equation}
    \mathcal{Y}_{train} \cap \mathcal{Y}_{val} = \varnothing, \quad
    \mathcal{Y}_{train} \cap \mathcal{Y}_{test} = \varnothing, \quad
    \mathcal{Y}_{val} \cap \mathcal{Y}_{test} = \varnothing
    \label{eq:disjoint_assert}
\end{equation}
Quy trình chỉ tiếp tục khi cả ba điều kiện thỏa mãn. Ngoài ra, một bộ kiểm thử tự động chạy trên dữ liệu thật xác nhận thứ tự thời gian train $<$ validation $<$ test, người dùng/sản phẩm của tập xác thực là tập con của tập huấn luyện, và người dùng/sản phẩm của tập kiểm thử là tập con của tập huấn luyện $\cup$ xác thực. Bảng~\ref{tab:split_hm} tổng hợp kết quả phân chia theo giao thức chính.

\begin{table}[htbp]
    \centering
    \caption{Kết quả phân chia theo mốc thời gian chung trên hm500k ($k=10$)}
    \label{tab:split_hm}
    \resizebox{\textwidth}{!}{%
    \begin{tabular}{|l|c|c|c|c|}
        \hline
        \textbf{Tập} & \textbf{Khoảng thời gian} & \textbf{Số cặp} & \textbf{Số người dùng} & \textbf{Số sản phẩm} \\
        \hline
        Huấn luyện & 20/09/2018 -- 30/06/2020 & 201.801 & 7.506 & 10.145 \\
        \hline
        Xác thực & 01/07/2020 -- 28/07/2020 & 7.057 & 2.275 & 2.642 \\
        \hline
        Kiểm thử & 29/07/2020 -- 22/09/2020 & 9.229 & 2.996 & 2.772 \\
        \hline
        Huấn luyện $\cup$ xác thực$^{*}$ & 20/09/2018 -- 28/07/2020 & 209.212 & 7.515 & 10.216 \\
        \hline
        \multicolumn{5}{|l|}{Giao nhau Train $\cap$ Val $=$ Train $\cap$ Test $=$ Val $\cap$ Test $= 0$ cặp} \\
        \hline
    \end{tabular}}
    \par\smallskip{\footnotesize $^{*}$Dữ liệu huấn luyện lại trước khi chấm kiểm thử: mọi cặp trước 29/07/2020, gồm cả 354 cặp của người dùng/sản phẩm mới xuất hiện trong giai đoạn xác thực (không được chấm ở tập xác thực vì chưa có trong tập huấn luyện).}
\end{table}

Mỗi người dùng trong tập huấn luyện có trung bình 26,9 cặp tương tác (trung vị 20); mỗi người dùng kiểm thử có trung bình 3,08 sản phẩm đúng. Khi chấm validation, tập ứng viên là 10.145 sản phẩm của tập huấn luyện, trừ những sản phẩm người dùng đã mua trong tập huấn luyện; khi chấm test, tập ứng viên là 10.216 sản phẩm của tập huấn luyện $\cup$ xác thực, trừ những sản phẩm người dùng đã mua trong tập này.

\section{Thiết kế kiến trúc mô hình NeuMF}
\label{sec:model_design}

\subsection{Cấu hình không gian nhúng Embedding độc lập}
\label{subsec:embedding_design}

Một quyết định thiết kế quan trọng của NeuMF là \emph{không} dùng chung một bảng Embedding cho hai nhánh GMF và MLP. Nhánh GMF tối ưu vector Embedding cho phép nhân Hadamard (tuyến tính, đòi hỏi các chiều đặc trưng ``thẳng hàng'' về ý nghĩa), còn nhánh MLP tối ưu vector Embedding làm đầu vào cho một mạng phi tuyến sâu. Ép hai mục tiêu khác bản chất này chia sẻ cùng một không gian tham số có thể khiến gradient từ hai nhánh xung đột nhau \parencite{he2017ncf}. Do đó mô hình dùng hai không gian nhúng độc lập:
\begin{itemize}
    \item Bảng Embedding nhánh GMF: $P^G \in \mathbb{R}^{M \times d_G}$ và $Q^G \in \mathbb{R}^{N \times d_G}$.
    \item Bảng Embedding nhánh MLP: $P^M \in \mathbb{R}^{M \times d_M}$ và $Q^M \in \mathbb{R}^{N \times d_M}$.
\end{itemize}
Việc tách hai bảng còn cho phép hai nhánh có số chiều khác nhau ($d_G \ne d_M$), như He và cộng sự cho phép trong công trình gốc -- điều cần thiết khi GMF và MLP tốt nhất sau tinh chỉnh có kích thước khác nhau (mục~\ref{sec:protocol}). Như công trình gốc, bảng Embedding được khởi tạo theo phân phối chuẩn có trung bình 0 và độ lệch chuẩn 0,01, tức $\mathcal{N}(0;\, 0{,}01^2)$; các tầng tuyến tính khởi tạo theo Xavier; lớp đầu ra $h$ không có hệ số chệch (bias), đúng với công thức $\sigma(h^T\phi)$. Số tham số chủ yếu nằm ở Embedding, tỉ lệ với $(M+N)\times d$: với cấu hình tốt nhất, GMF ($d=64$) có 1.143.360 tham số, MLP ($d=32$) có 574.400, NeuMF-Scratch ($d=64$) có 2.297.536 và NeuMF-Pretrained ($d_G=64$, $d_M=32$) có 1.717.760; BPR-MF ($d=32$) có 571.648.

\subsection{Thiết kế nhánh Generalized Matrix Factorization (GMF)}
\label{subsec:gmf_design}

Nhánh GMF là dạng \emph{tổng quát hóa} của nhân tử hóa ma trận cổ điển (Công thức~\eqref{eq:mf_dot}): thay vì cộng dồn ngay tích các chiều thành một số vô hướng ($p_u^T q_i$), GMF giữ kết quả nhân theo từng chiều dưới dạng vector qua phép nhân Hadamard:
\begin{equation}
    \phi^{GMF} = p_u^G \odot q_i^G
\end{equation}
Mỗi chiều $f$ của $\phi^{GMF} \in \mathbb{R}^{d_G}$ đo mức độ ``khớp'' giữa User $u$ và Item $i$ theo đặc trưng ẩn thứ $f$. Phép cộng dồn được chuyển sang lớp trọng số $h$ ở tầng đầu ra (mục~\ref{subsec:neumf_output}), cho phép mô hình học \emph{trọng số riêng cho từng chiều} thay vì coi mọi chiều quan trọng như nhau -- đây là phần ``tổng quát hóa'' trong tên gọi GMF; khi $h = \mathbf{1}$, GMF trở về MF.

\subsection{Thiết kế nhánh Multi-Layer Perceptron (MLP)}
\label{subsec:mlp_design}

Nhánh MLP \emph{nối} (concatenate) hai vector Embedding thành một vector đầu vào duy nhất, không áp đặt trước cách hai vector tương tác -- việc ``học cách tương tác'' được giao cho các tầng mạng phía sau:
\begin{equation}
    z_1 = \phi_1^{MLP} = \begin{bmatrix} p_u^M \\ q_i^M \end{bmatrix} \in \mathbb{R}^{2 d_M}
\end{equation}
Mạng MLP có kiến trúc tháp giảm một nửa qua mỗi tầng ẩn, theo quy ước của NCF gốc:
\begin{equation}
    [2d_M \longrightarrow d_M \longrightarrow d_M/2 \longrightarrow d_M/4] \qquad (d_M = 32:\; [64 \to 32 \to 16 \to 8])
\end{equation}
Mỗi tầng ẩn gồm một phép biến đổi tuyến tính, hàm kích hoạt ReLU và lớp Dropout (tỉ lệ được tinh chỉnh trong $\{0;\, 0{,}2\}$) để hạn chế quá khớp -- rủi ro thực sự khi dữ liệu huấn luyện chỉ có 201.801 cặp dương.

\subsection{Nhánh MLP đứng riêng: mô hình DNN thuần, không phải mô hình hợp nhất}
\label{subsec:early_fusion_design}

Khi đứng riêng, nhánh MLP nối Embedding của User và Item ngay ở đầu vào rồi đưa qua một tháp nơ-ron:
\begin{equation}
    z_1 = \begin{bmatrix} p_u \\ q_i \end{bmatrix} \in \mathbb{R}^{2d}, \qquad
    \hat{y}_{ui}^{MLP} = \sigma\left(h^T \cdot \text{MLP}(z_1)\right)
    \label{eq:mlp_standalone}
\end{equation}
Mô hình này chỉ có view phi tuyến, không có thành phần MF nào, nên trong đề tài nó là \textbf{đối chứng DNN thuần} (tương đương bỏ nhánh MF khỏi NeuMF), không phải một mô hình hợp nhất theo nghĩa ở mục~\ref{subsec:early_late_fusion}. Trong mã nguồn, nhóm từng cài đặt riêng một lớp \texttt{EarlyFusionModel} theo cách hiểu ``nối biểu diễn User--Item ngay ở đầu vào''; đối chiếu với mục~\ref{subsec:mlp_design} cho thấy lớp này \emph{trùng hoàn toàn} với MLP đứng riêng (cùng cặp bảng Embedding, cùng phép nối, cùng tháp và lớp đầu ra). Điều này được kiểm chứng bằng thực nghiệm (NDCG@10 trên validation giống hệt nhau với cùng seed) và bằng một bài kiểm thử tự động so sánh kiến trúc, nên lớp này không được tinh chỉnh hay đánh giá riêng (đã ghi trong kế hoạch thực nghiệm trước khi chấm test, mục~\ref{sec:protocol}). Phép hợp nhất MF với DNN của đề tài nằm ở tầng kết hợp của NeuMF (mục~\ref{subsec:neumf_output}).

\subsection{Tầng kết hợp NeuMF và lớp dự đoán đầu ra}
\label{subsec:neumf_output}

Đây là bước hợp nhất của mô hình lai -- Early Fusion ở mức biểu diễn theo phân loại ở mục~\ref{subsec:early_late_fusion}: vector của nhánh GMF ($\phi^{GMF}$, tương tác tuyến tính) và của nhánh MLP ($\phi_L^{MLP}$, tương tác phi tuyến) được ghép nối thành một vector hợp nhất:
\begin{equation}
    \phi^{NeuMF} = \begin{bmatrix} p_u^G \odot q_i^G \\ \phi_L^{MLP} \end{bmatrix}
\end{equation}
Vector này qua một tầng tuyến tính cuối với trọng số $h = [h_G^T, h_M^T]^T$ và hàm Sigmoid:
\begin{equation}
    \hat{y}_{ui} = \sigma \left( h^T \phi^{NeuMF} \right) = \sigma \left( h_G^T (p_u^G \odot q_i^G) + h_M^T \phi_L^{MLP} \right)
\end{equation}
Điểm dự đoán là một tổ hợp tuyến tính có trọng số học được giữa ``quan điểm'' của nhánh tuyến tính và nhánh phi tuyến. Với NeuMF-Pretrained, lúc khởi tạo $h = [\alpha\, h_{GMF};\, (1-\alpha)\, h_{MLP}]$ nên logit ban đầu đúng bằng $\alpha$ lần logit của GMF cộng $(1-\alpha)$ lần logit của MLP -- giống một phép trộn điểm; điểm khác với Late Fusion là sau đó toàn bộ mạng, kể cả hai nhánh, được tinh chỉnh chung thay vì giữ nguyên hai mô hình thành phần. Khi xếp hạng, hệ thống dùng trực tiếp giá trị trước Sigmoid (logit): Sigmoid là hàm đơn điệu nên không đổi thứ tự, nhưng ở độ chính xác float32 nó bão hòa về đúng 1,0 với điểm lớn và tạo ra các giá trị bằng nhau giả.

\section{Kỹ thuật lấy mẫu tiêu cực và hàm mục tiêu huấn luyện}
\label{sec:sampling_loss}

\subsection{Kỹ thuật lấy mẫu tiêu cực động (Dynamic Negative Sampling)}
\label{subsec:dynamic_sampling}

Dữ liệu phản hồi ẩn (mục~\ref{subsec:feedback_types}) chỉ ghi nhận tương tác đã xảy ra ($y_{ui}=1$), không có nhãn âm tường minh. Negative Sampling giả định các sản phẩm người dùng chưa tương tác nhiều khả năng là âm (dù một phần trong đó chỉ là người dùng chưa biết đến). Cụ thể, trong mỗi epoch, với mỗi cặp dương $(u, i) \in \mathcal{Y}_{train}$, hệ thống lấy ngẫu nhiên $N_{ratio}$ sản phẩm $j \in \mathcal{I} \setminus \mathcal{I}_u^{train}$ và gán $y_{uj} = 0$. Tập loại trừ chỉ gồm sản phẩm của người dùng trong \emph{tập huấn luyện}, để không dùng nhãn tương lai của tập xác thực/kiểm thử. $N_{ratio}$ được tinh chỉnh trong $\{4, 8\}$: He và cộng sự \parencite{he2017ncf} ghi nhận 1 mẫu âm là không đủ và tỉ lệ khoảng 3--6 cho kết quả tốt nhất. Mẫu âm được lấy lại hoàn toàn ở \emph{mỗi} epoch (lấy mẫu động), giúp mô hình thấy nhiều tổ hợp cặp âm khác nhau và giảm nguy cơ học vẹt một tập âm cố định.

\subsection{Hàm tổn thất Binary Cross-Entropy (BCE) và Regularization}
\label{subsec:bce_regularization}

Với cả nhãn dương và nhãn âm, huấn luyện trở thành bài toán phân loại nhị phân. Hàm Binary Cross-Entropy được tối thiểu hóa, cùng số hạng điều chuẩn $L_2$ (weight decay của bộ tối ưu):
\begin{equation}
    \mathcal{L} = -\sum_{(u, i) \in \mathcal{Y} \cup \mathcal{Y}^-} \left[ y_{ui} \log \hat{y}_{ui} + (1 - y_{ui}) \log (1 - \hat{y}_{ui}) \right] + \lambda_\Theta \|\Theta\|_2^2
    \label{eq:bce_loss_full}
\end{equation}
trong đó $\Theta$ là toàn bộ tham số có thể học và $\lambda_\Theta \in \{0;\, 10^{-6}\}$ được tinh chỉnh. Trong cài đặt, BCE được tính trên logit (\texttt{BCEWithLogitsLoss}) để ổn định số học. Cực tiểu hóa $\mathcal{L}$ tương đương với huấn luyện mô hình phân biệt tốt nhất giữa cặp thật và cặp lấy mẫu ngẫu nhiên.

\section{Quy trình huấn luyện và tối ưu hóa mô hình}
\label{sec:training_pipeline}

\subsection{Chiến lược tiền huấn luyện (Pre-training GMF + MLP)}
\label{subsec:pretraining_strategy}

Hàm mất mát của NeuMF không lồi và không gian tham số lớn, nên khởi tạo ngẫu nhiên toàn bộ mạng (NeuMF-Scratch) có thể hội tụ về một cực tiểu kém. Chiến lược tiền huấn luyện của công trình gốc \parencite{he2017ncf} chia quá trình tối ưu thành hai giai đoạn:
\begin{enumerate}
    \item \textbf{Huấn luyện trước:} Huấn luyện độc lập GMF và MLP bằng Adam \parencite{kingma2014adam} cho tới khi dừng sớm, lưu Embedding và trọng số tốt nhất của từng nhánh.
    \item \textbf{Hợp nhất và tinh chỉnh:} Khởi tạo bốn bảng Embedding và các tầng MLP của NeuMF từ hai mô hình trên; khởi tạo trọng số đầu ra $h \leftarrow [\alpha\, h_{GMF}^T, (1-\alpha)\, h_{MLP}^T]^T$, trong đó $\alpha$ cân bằng đóng góp ban đầu của hai nhánh (bài gốc dùng $\alpha=0{,}5$; dự án tinh chỉnh $\alpha \in \{0{,}3;\, 0{,}5;\, 0{,}7\}$); sau đó tinh chỉnh toàn mạng với tốc độ học nhỏ (tinh chỉnh trong $\{10^{-3}, 5\cdot10^{-4}, 10^{-4}\}$).
\end{enumerate}
Khác với bài gốc (tinh chỉnh bằng SGD), dự án dùng \textbf{Adam cho cả hai giai đoạn}: trên dữ liệu này, SGD với tốc độ học 0,01 khiến NeuMF không hội tụ (hàm mất mát dừng ở mức entropy của nhãn). NeuMF-Pretrained và NeuMF-Scratch được huấn luyện và đánh giá song song trong mọi thực nghiệm để tách riêng đóng góp của Pre-training.

\subsection{Bảng thiết lập siêu tham số và kỹ thuật dừng sớm (Early Stopping)}
\label{subsec:hyperparameters_setup}

Bảng~\ref{tab:hyperparams} liệt kê cấu hình mặc định và không gian tinh chỉnh. Cấu hình mặc định luôn là một trong các cấu hình được thử; cấu hình dùng cho kết quả cuối của từng mô hình là cấu hình tốt nhất trên tập xác thực (Bảng~\ref{tab:tuning}).

\begin{table}[htbp]
    \centering
    \caption{Siêu tham số: giá trị mặc định và không gian tinh chỉnh}
    \label{tab:hyperparams}
    \begin{tabularx}{\textwidth}{|>{\raggedright\arraybackslash}X|c|c|}
        \hline
        \textbf{Siêu tham số} & \textbf{Mặc định} & \textbf{Không gian tinh chỉnh} \\
        \hline
        Kích thước lô (Batch Size) & $512$ & cố định \\
        \hline
        Bộ tối ưu & Adam & cố định \\
        \hline
        Tốc độ học -- GMF, MLP, NeuMF-Scratch & $10^{-3}$ & $\{10^{-3},\; 5\cdot 10^{-4}\}$ \\
        \hline
        Tốc độ học tinh chỉnh -- NeuMF-Pretrained & $10^{-3}$ & $\{10^{-3},\; 5\cdot 10^{-4},\; 10^{-4}\}$ \\
        \hline
        Hệ số $\alpha$ khởi tạo lớp đầu ra (Pretrained) & $0{,}5$ & $\{0{,}3;\; 0{,}5;\; 0{,}7\}$ \\
        \hline
        Kích thước Embedding $d$ & $32$ & $\{32, 64\}$ \\
        \hline
        Cấu trúc tầng ẩn MLP & $[64, 32, 16, 8]$ & $[2d, d, d/2, d/4]$ \\
        \hline
        Tỷ lệ lấy mẫu âm $N_{ratio}$ & $4$ & $\{4, 8\}$ \\
        \hline
        Tỷ lệ Dropout (MLP, NeuMF) & $0{,}2$ & $\{0;\; 0{,}2\}$ \\
        \hline
        Weight Decay $\lambda$ & $10^{-6}$ & $\{0;\; 10^{-6}\}$ \\
        \hline
        Số epoch tối đa & $20$ & cố định \\
        \hline
        Early Stopping: theo dõi / kiên nhẫn & NDCG@10 val / 5 & cố định \\
        \hline
        BPR-MF: kích thước $d$ & $32$ & $\{32, 64, 128\}$ \\
        \hline
        BPR-MF: tốc độ học & $0{,}03$ & $\{0{,}01;\; 0{,}03;\; 0{,}05\}$ \\
        \hline
        BPR-MF: hệ số điều chuẩn & $0{,}005$ & $\{0{,}001;\; 0{,}005;\; 0{,}01\}$ \\
        \hline
        BPR-MF: số epoch & $20$ & $\{20, 40\}$ \\
        \hline
    \end{tabularx}
\end{table}

Early Stopping: sau mỗi epoch, mô hình được chấm NDCG@10 theo Full Ranking trên tập xác thực; trạng thái tốt nhất được giữ lại, và huấn luyện dừng khi chỉ số không cải thiện sau 5 epoch liên tiếp.

\section{Quy trình tinh chỉnh và đánh giá không thiên lệch}
\label{sec:protocol}

Ferrari Dacrema và cộng sự \parencite{dacrema2019progress} chỉ ra mã nguồn gốc của NCF chọn số epoch dựa trên chính tập kiểm thử, và các baseline trong nhiều bài báo học sâu được tinh chỉnh sơ sài; Rendle và cộng sự \parencite{rendle2020ncfmf} cho thấy MF được tune kỹ vượt NeuMF trên chính dữ liệu của bài gốc. Để kết luận của đề tài không mắc các lỗi này, quy trình thực nghiệm áp dụng năm nguyên tắc:
\begin{enumerate}
    \item \textbf{Khoá tập kiểm thử.} Mọi lựa chọn mô hình và siêu tham số chỉ dựa trên tập xác thực. Tập kiểm thử chỉ được chấm qua kịch bản đánh giá cuối (và kịch bản phân tích phụ, vốn chỉ nạp lại đúng các trọng số của kịch bản đó); các kịch bản này từ chối chạy nếu kế hoạch thực nghiệm chưa được lưu vào hệ thống quản lý phiên bản, nếu thiếu lý do chạy, hoặc nếu mã nguồn có thay đổi chưa lưu. Mỗi lần chấm test ghi một dòng vào nhật ký truy cập (thời gian, phiên bản mã nguồn, mã băm cấu hình, lý do).
    \item \textbf{Đăng ký trước (pre-registration).} Độ đo chính, lưới tinh chỉnh, số seed, họ so sánh và tiêu chí kết luận được cố định và lưu lại \emph{trước} lần chấm test đầu tiên. Mọi thay đổi sau đó được ghi vào mục ``lệch kế hoạch'' kèm thời điểm (trước hay sau khi xem test).
    \item \textbf{Cùng ngân sách tinh chỉnh.} Mỗi mô hình có siêu tham số được thử đúng \textbf{6 cấu hình}, lấy bằng một lần rút ngẫu nhiên cố định từ lưới của mô hình đó (Bảng~\ref{tab:hyperparams}), cấu hình mặc định luôn nằm trong 6. Mỗi cấu hình ghi một dòng nhật ký (mã băm cấu hình, phiên bản mã nguồn, chỉ số trên validation, thời gian). NeuMF-Pretrained được nạp từ GMF và MLP tốt nhất.
    \item \textbf{Chọn mô hình lai trên validation.} Trong hai biến thể NeuMF, mô hình có NDCG@10 validation cao nhất được xem là ``mô hình lai được chọn''; tập test chỉ để báo cáo, không để chọn lại.
    \item \textbf{Kiểm định thống kê và tiêu chí kết luận cố định.} Mỗi mô hình (với cấu hình tốt nhất) được chạy với 3 seed (42, 2024, 2025); với mỗi seed, mô hình được huấn luyện trên tập huấn luyện để chọn số epoch theo validation, rồi huấn luyện lại trên tập huấn luyện $\cup$ xác thực đúng số epoch đó và chấm trên test (mục~\ref{subsec:split_sampling}). NDCG@10 của từng người dùng được lấy trung bình qua 3 seed, rồi so sánh từng cặp mô hình bằng kiểm định Wilcoxon signed-rank theo người dùng \parencite{wilcoxon1945} và khoảng tin cậy 95\% bằng paired bootstrap (10.000 lần lấy mẫu lại) \parencite{efron1993bootstrap}; giá trị $p$ được hiệu chỉnh Holm \parencite{holm1979} trên họ 8 so sánh cố định (NeuMF-Pretrained so với BPR-MF, Most Popular, GMF, MLP, NeuMF-Scratch; NeuMF-Scratch so với BPR-MF, GMF, MLP). Chỉ viết ``A tốt hơn B'' khi đồng thời: $p_{Holm} < 0{,}05$, khoảng tin cậy không chứa 0, và chênh lệch tương đối NDCG@10 $\ge 5\%$; ngược lại kết luận ``không khác biệt có ý nghĩa''. Không dùng Wilcoxon theo seed vì với 3 seed, giá trị $p$ nhỏ nhất có thể đạt là $2/2^3 = 0{,}25$.
\end{enumerate}

\paragraph{Lệch kế hoạch trước lần chấm test đầu tiên.} Kế hoạch ban đầu có thêm các ứng viên iALS \parencite{rendle2022ials}, Most Popular theo cửa sổ thời gian gần, DeepCF/CFNet và một mô hình trộn điểm iALS + NeuMF-Scratch (trộn một mô hình MF với mô hình lai NeuMF); tất cả đã được tinh chỉnh trên validation nhưng sau đó bị loại khỏi phạm vi theo quyết định của nhóm (đã quá nhiều baseline), \emph{trước} khi xem bất kỳ số liệu test nào. Kết quả validation của chúng vẫn được giữ trong nhật ký tinh chỉnh để minh bạch và được nhắc lại ở mục~\ref{sec:discussion}. Lớp \texttt{EarlyFusionModel} (nối Embedding ngay ở đầu vào) bị loại vì trùng kiến trúc MLP (mục~\ref{subsec:early_fusion_design}). Bản ItemKNN đầu tiên dùng ma trận tương đồng dày $N \times N$ và bị chặn ở 5.000 sản phẩm, nên không chạy được với $N = 10.345$; bản thưa top-$k$ được bổ sung ở phần mở rộng dưới đây.

\paragraph{Lệch kế hoạch sau lần chấm test đầu tiên.} Lần chấm đầu (27/09/2026) huấn luyện các mô hình chỉ trên tập huấn luyện rồi chấm thẳng tập kiểm thử, bỏ trống 4 tuần của giai đoạn xác thực. Sau đó nhóm bổ sung bước huấn luyện lại trên tập huấn luyện $\cup$ xác thực (mục~\ref{subsec:split_sampling}) và chấm lại toàn bộ (29/09/2026); cấu hình, độ đo, họ so sánh và tiêu chí kết luận giữ nguyên. Vì thay đổi này được quyết định \emph{sau} khi đã xem số liệu test, số liệu ở Chương~\ref{chap:thucnghiem} -- lấy từ lần chấm thứ hai -- cần được hiểu là kết quả sau khi đã xem test một lần. Kết luận không đổi giữa hai lần chấm: ở lần đầu, BPR-MF cũng có NDCG@10 trung bình cao nhất (0,01062, so với 0,01001 của NeuMF-Pretrained và 0,00958 của NeuMF-Scratch) và cũng 0/8 so sánh đạt tiêu chí ``tốt hơn''. Nhật ký truy cập ghi đủ cả hai lần chấm; số liệu của lần đầu được giữ trong lịch sử phiên bản của kho mã.

\paragraph{Mở rộng sau khi xem test.} Để trả lời câu hỏi Early Fusion hay Late Fusion phù hợp hơn và bổ sung baseline láng giềng mà đề cương đặt ra, đề tài bổ sung bốn mô hình: Late Fusion GMF + MLP và Late Fusion BPR-MF + MLP (trộn điểm đã chuẩn hóa của hai mô hình huấn luyện riêng, công thức~\eqref{eq:late_fusion_mfdnn}), ItemKNN và UserKNN (láng giềng top-$k$ theo cosine, mục~\ref{subsec:memory_model_based}). Tham số được chọn chỉ trên tập xác thực: ItemKNN và UserKNN thử 6 cấu hình như mọi mô hình; Late Fusion quét đủ 11 giá trị $w \in \{0;\, 0{,}1;\, \dots;\, 1\}$ (một tham số, không huấn luyện thêm mạng nào). Lựa chọn được ghi vào mục ``mở rộng'' của kế hoạch đăng ký trước \emph{trước} khi chấm test, kèm một lần chạy lại kiểm tra cho đúng cùng số trên cả 34 cấu hình. Khi chấm test, Late Fusion dùng đúng các mô hình GMF, MLP, BPR-MF đã huấn luyện lại trên tập huấn luyện $\cup$ xác thực của từng seed (không huấn luyện thêm), còn ItemKNN, UserKNN tính trên tập huấn luyện $\cup$ xác thực. Kiểm định dùng một họ 7 so sánh riêng có hiệu chỉnh Holm riêng (Late Fusion GMF + MLP so với NeuMF-Scratch, NeuMF-Pretrained, GMF, MLP; Late Fusion BPR-MF + MLP so với BPR-MF; ItemKNN và UserKNN so với NeuMF-Scratch), nên không ảnh hưởng tới họ 8 so sánh chính. Vì phần này được thêm sau khi đã xem số liệu test, kết quả của nó được báo cáo là kết quả mở rộng, không thay thế kết luận chính.

\paragraph{Chạy lại toàn bộ trên phiên bản mã nguồn cuối.} Kết quả lần chấm thứ hai được sinh ở một phiên bản mã nguồn cũ hơn; sau đó thư mục mã được đổi tên và phần mở rộng được thêm vào. Để mọi số liệu trong báo cáo truy vết về \emph{cùng một} phiên bản, toàn bộ chuỗi đánh giá cuối được chạy lại một lần trên phiên bản cuối (đăng ký trước trong kế hoạch, cùng cấu hình, seed, độ đo, họ so sánh và tiêu chí; không tinh chỉnh lại). Vì huấn luyện chạy ở chế độ GPU tất định và mã huấn luyện, đánh giá của bảy mô hình chính không đổi, lần chạy lại trên cùng máy được kỳ vọng ra đúng cùng số -- đây đồng thời là một kiểm tra tái lập; kết quả đối chiếu được ghi ở mục~\ref{sec:exp_config}. Nhật ký truy cập ghi đủ mọi lần chấm.

\section{Thiết kế hệ thống giao diện thực nghiệm (Demo)}
\label{sec:demo_design}

Mục này trình bày hệ thống demo minh họa mô hình đã huấn luyện trong một kịch bản cửa hàng thời trang. Demo tách biệt về kiến trúc với phần huấn luyện: chỉ nạp trọng số mô hình đã lưu và gọi lại đúng các hàm tiền xử lý, xếp hạng của pipeline huấn luyện, không huấn luyện lại.

\subsection{Kiến trúc tổng thể: ba tầng tách biệt}
\label{subsec:demo_architecture}

\begin{figure}[htbp]
    \centering
    \resizebox{0.92\textwidth}{!}{%
    \begin{tikzpicture}[
        node distance=1.0cm,
        stage/.style={rectangle, draw=black!60, very thick, rounded corners=2pt, text width=11cm, align=left, font=\small, inner sep=8pt},
        arrow/.style={->, >=stealth, thick, color=black!70}
    ]
        \node[stage, fill=gray!12] (client) at (0,0)
            {\textbf{TẦNG 1 -- GIAO DIỆN NGƯỜI DÙNG (Frontend)}\\[2pt]
             Trang HTML/JavaScript tĩnh, biểu đồ bằng Chart.js. Ba màn hình: Khách hàng mới, Admin kiểm thử mô hình, Dashboard. Chỉ gọi API, không truy cập trực tiếp trọng số hay dữ liệu.};

        \node[stage, fill=blue!8, below=of client] (api)
            {\textbf{TẦNG 2 -- API SUY DIỄN (FastAPI)}\\[2pt]
             Các endpoint ở Bảng~\ref{tab:demo_api}: kiểm tra đầu vào, tra cứu ID, gọi tầng mô hình, định dạng kết quả. Suy diễn thuần (inference-only), không lan truyền ngược.};

        \node[stage, fill=green!8, below=of api] (model)
            {\textbf{TẦNG 3 -- MÔ HÌNH \& DỮ LIỆU TRA CỨU}\\[2pt]
             Dựng lại dữ liệu theo đúng cấu hình của đánh giá cuối, đối chiếu với tệp kết quả đã lưu (số người dùng kiểm thử, số sản phẩm ứng viên) rồi mới nạp trọng số GMF/MLP/NeuMF và BPR-MF với cấu hình riêng của từng mô hình; dựng ItemKNN/UserKNN và Late Fusion của phần mở rộng; bảng ánh xạ ID gốc $\leftrightarrow$ mã nội bộ.};

        \draw[arrow] (client) -- node[right, font=\scriptsize] {HTTP/JSON} (api);
        \draw[arrow] (api) -- node[right, font=\scriptsize] {nạp 1 lần, lưu cache} (model);
    \end{tikzpicture}}
    \caption{Kiến trúc ba tầng của hệ thống giao diện thực nghiệm (demo)}
    \label{fig:demo_architecture}
\end{figure}

Hai nguyên tắc thiết kế quan trọng nhất: \textit{(i)} tách suy diễn khỏi huấn luyện, nên thời gian phản hồi không phụ thuộc chi phí huấn luyện và tầng 2+3 có thể đóng gói thành một container độc lập; \textit{(ii)} demo \emph{tự kiểm tra} rằng dữ liệu dựng lại khớp với lần chạy đã huấn luyện -- nếu lệch (ví dụ sai tệp dữ liệu hoặc sai ngưỡng $k$), máy chủ báo lỗi thay vì nạp trọng số với ánh xạ ID sai. Bộ kiểm thử tự động của demo xác nhận tập ứng viên trùng với tập ứng viên của đánh giá cuối và các chỉ số theo người dùng mà demo tính lại khớp tệp kết quả theo người dùng của Chương~\ref{chap:thucnghiem}; chỉ vài cặp sản phẩm gần như hoà điểm ở hạng rất sâu (trên 100) có thể đổi chỗ vì logit được tính lại trên CPU thay vì GPU, không ảnh hưởng chỉ số @K.

\subsection{Thiết kế API suy diễn}
\label{subsec:demo_api}

\begin{table}[htbp]
    \centering
    \caption{Các endpoint chính của API suy diễn}
    \label{tab:demo_api}
    \begin{tabularx}{\textwidth}{|l|X|}
        \hline
        \textbf{Endpoint} & \textbf{Chức năng} \\
        \hline
        \texttt{GET /api/context} & Thông tin bộ dữ liệu, lần chạy đang dùng, số người dùng/sản phẩm, các $K$ hợp lệ, mô hình khả dụng. \\
        \hline
        \texttt{GET /api/users/search}, \texttt{/random} & Tìm hoặc chọn ngẫu nhiên một khách hàng có sản phẩm kiểm thử, lọc theo nhóm ít/trung bình/nhiều giao dịch. \\
        \hline
        \texttt{GET /api/users/\{id\}/history} & Lịch sử mua trong tập huấn luyện, sắp theo thời gian. \\
        \hline
        \texttt{GET /api/users/\{id\}/recommend} & Top-$K$ của một mô hình, hạng của từng sản phẩm đích, số ứng viên, HR@$K$, NDCG@$K$, Recall@$K$ của khách đó. \\
        \hline
        \texttt{GET /api/users/\{id\}/neighbors} & Top-$K$ khách tương đồng nhất theo UserKNN: độ tương đồng, số sản phẩm mua chung, láng giềng nào đã mua sản phẩm đích của khách. \\
        \hline
        \texttt{POST /api/onboarding/recommend} & Gợi ý cho khách hàng mới theo luật (lọc theo lựa chọn + độ phổ biến trong train) -- được gắn nhãn rõ \emph{không phải mô hình NeuMF}. \\
        \hline
        \texttt{GET /api/dashboard} & Số liệu tổng hợp đọc từ tệp kết quả đã chạy, kèm tên tệp nguồn. \\
        \hline
        \texttt{POST /api/reload} & Xoá cache, nạp lần chạy mới nhất mà không cần khởi động lại. \\
        \hline
    \end{tabularx}
\end{table}

Luồng xử lý của chức năng gợi ý Top-$K$ gồm 5 bước, dùng lại đúng các hàm của pipeline đánh giá:
\begin{enumerate}
    \item Tra định danh khách hàng để lấy mã nội bộ. Khách chưa có trong dữ liệu huấn luyện là tình huống Cold-start tuyệt đối, nằm ngoài khả năng của mô hình (mục~\ref{sec:pham_vi_thuc_nghiem}); demo chuyển sang màn hình Khách hàng mới với gợi ý theo luật thay vì suy diễn một kết quả không đáng tin cậy.
    \item Xác định tập ứng viên: toàn bộ sản phẩm có trong dữ liệu huấn luyện trừ những sản phẩm khách đã mua -- cùng quy tắc dựng tập ứng viên của đánh giá Full Ranking.
    \item Tính điểm cho toàn bộ tập ứng viên theo lô (vector hoá), dùng logit với mô hình nơ-ron.
    \item Sắp xếp giảm dần với cùng quy tắc phá hòa tất định như lúc đánh giá, lấy $K$ phần tử đầu.
    \item Ánh xạ ngược mã nội bộ về \texttt{article\_id}, kèm tên và loại sản phẩm, trả về Tầng 1.
\end{enumerate}

\subsection{Lựa chọn mô hình triển khai và yêu cầu phi chức năng}
\label{subsec:demo_model_choice}

Demo cho phép chọn và so sánh song song hai mô hình bất kỳ trong số NeuMF-Pretrained, NeuMF-Scratch, GMF, MLP, BPR-MF, Most Popular và bốn mô hình của phần mở rộng, để người xem tự đối chiếu với kết luận của Chương~\ref{chap:thucnghiem}; demo không khẳng định mô hình nào ``tốt nhất'' ngoài những gì kiểm định thống kê cho phép. Demo nạp đúng các trọng số của đánh giá cuối (seed 42): mô hình đã huấn luyện lại trên tập huấn luyện $\cup$ xác thực, lịch sử hiển thị là đúng dữ liệu mô hình đã học, và sản phẩm đích là các món khách mua lần đầu trong giai đoạn kiểm thử -- nên số liệu trên demo là hiển thị lại kết quả của Chương~\ref{chap:thucnghiem}, không phải một lần chấm mới. Màn hình Admin còn hiện 10 khách tương đồng nhất theo UserKNN để minh họa trực tiếp cách lọc cộng tác dựa trên láng giềng chọn ra gợi ý (mục~\ref{subsec:memory_model_based}). Chế độ khám phá cũ (lần chạy Leave-One-Out với cấu hình mặc định) vẫn được giữ để đối chiếu. Về yêu cầu phi chức năng, mỗi yêu cầu gợi ý chỉ là một lượt suy diễn thuận trên khoảng 10.000 ứng viên, nên mục tiêu thiết kế là độ trễ mức mili giây trên CPU; độ trễ thực đo (trung vị và phân vị 95) được báo cáo ở mục~\ref{subsec:latency}.

\subsection{Minh họa hệ thống demo đã cài đặt}
\label{subsec:demo_screenshots}

Hệ thống đã được cài đặt và có bộ kiểm thử tự động riêng (kiểm tra khớp số liệu với đánh giá cuối, đúng quy tắc tập ứng viên, đúng định dạng API). Hình~\ref{fig:demo_screenshot_hm} minh họa màn hình Admin kiểm thử mô hình và Hình~\ref{fig:demo_neighbors} là bảng khách tương đồng của cùng khách hàng.

\hinhketqua{media/figures/demo/demo_final.png}{\textwidth}{Giao diện demo -- màn hình Admin kiểm thử mô hình với mô hình của đánh giá cuối: lịch sử mua, Top-10 của NeuMF-Pretrained và BPR-MF, hạng của từng sản phẩm đích trong tập kiểm thử}{fig:demo_screenshot_hm}{chụp màn hình demo (\texttt{python run.py demo}) và lưu thành \texttt{media/figures/demo/demo\_final.png}}

\hinhketqua{media/figures/demo/demo_neighbors.png}{\textwidth}{Giao diện demo -- 10 khách hàng tương đồng nhất theo UserKNN: độ tương đồng cosine, số sản phẩm mua chung và láng giềng đã mua sản phẩm đích nào của khách}{fig:demo_neighbors}{chụp màn hình demo và lưu thành \texttt{media/figures/demo/demo\_neighbors.png}}
